\chapter{设计}
\label{chap:design}

\sysname 将多播-聚合流实现为 channel，并以组播 credit 逐 packet offset
推进。Credit Context 绑定 offset 与 subchannel，credit 副本沿下行路径取得
逐级聚合位置；requester 带回这些位置后，request 在上行路径完成聚合。
responder 再以完成事件选择 subchannel、约束后续 credit，并在异常时启动
整组 repair。\Cref{tab:design-sync} 汇总这些规则如何满足
\cref{ssec:inc-limits} 的上下文一致性与整组进度要求。

\begin{table}[t]
    \centering
    \caption{\sysname 用同一组协议字段建立聚合上下文一致性并暴露整组进度。}
    \label{tab:design-sync}
    \small
    \setlength{\tabcolsep}{3pt}
    \begin{tabularx}{\linewidth}{p{0.18\linewidth}YY}
        \toprule
        机制 & 上下文绑定 & 整组进度 \\
        \midrule
        多播-聚合流 &
        定义路径、聚合位置与恢复的共同 packet-offset 边界 &
        把 responder 与 requester group 的 request/response 关系作为一个传输对象 \\
        组播 credit &
        显式绑定 offset、subchannel 与逐级 slot &
        requester group 消费同一份 credit，并继承同一 Credit Context \\
        聚合器位置反馈 &
        response 下行分配 slot，request 上行带回并通过 owner 校验 &
        同一 offset 的所有输入在其共享的每一级命中同一聚合位置 \\
        拥塞控制 &
        responder 以 $r_s$ 和 $W_s$ 两道门约束后续 credit，repair 使用独立预算 &
        fan-in 完成时间反映最慢分支，一份闭环完成时延样本约束整组速率 \\
        可靠性 &
        owner mismatch 生成失配通知，repair request 走旁路，二者都不产生正常完成样本 &
        responder 组播 repair trigger，并用 requester 位图完成整组去重 \\
        \bottomrule
    \end{tabularx}
\end{table}

\section{总体设计}
\label{ssec:design-overview}

运行时把一次原语调用映射为一个或多个 channel message。每个 message 覆盖
一段连续数据，并进一步切成固定大小的 packet。以 \texttt{AllReduce} 为例，
数据按 rank 切片，每个 rank 作为一个切片的 responder，其它 rank 为该切片
提供同一 packet offset 上的输入；其余原语的映射见
\cref{ssec:channel-subchannel}。

\begin{figure*}[t]
\centering
% 设计总览：多播-聚合流与组播 credit（fig:design-overview）
% (a) 一份 normal credit 的组播闭环；(b) 一个 channel 的多棵 subchannel 树
\begin{minipage}[b]{0.56\linewidth}
\centering
\begin{tikzpicture}[x=1cm,y=1cm,
  box/.style={draw,rounded corners=1.5pt,minimum height=0.5cm,
    inner sep=3pt,font=\scriptsize},
  acc/.style={box,fill=ArborLightGray!60},
  ep/.style={box,fill=ArborBlue!12},
  credit/.style={-stealth,ArborBlue,semithick,dashed},
  up/.style={-stealth,ArborOrange,semithick},
  lab/.style={font=\tiny,inner sep=1pt}]
  \node[ep]  (R)  at ( 0.0,4.55) {responder};
  \node[acc] (SP) at ( 0.0,3.15) {加速器};
  \node[acc] (L1) at (-1.9,1.75) {加速器};
  \node[acc] (L2) at ( 1.9,1.75) {加速器};
  \node[ep]  (q1) at (-2.7,0.35) {$q_1$};
  \node[ep]  (q2) at (-1.1,0.35) {$q_2$};
  \node[ep]  (q3) at ( 1.9,0.35) {$q_3$};
  \node[lab,anchor=west] at ( 0.62,3.38) {slot 7};
  \node[lab,anchor=east] at (-2.52,1.98) {slot 3};
  \node[lab,anchor=west] at ( 2.52,1.98) {slot 5};
  % 下行：credit 组播副本，逐级推进 allocator 并压栈
  \draw[credit] (R)  to[bend right=18]
    node[lab,left=2pt]{credit $(k,s_1)$} (SP);
  \draw[credit] (SP) to[bend right=18] (L1);
  \draw[credit] (SP) to[bend right=18] (L2);
  \draw[credit] (L1) to[bend right=18]
    node[lab,left=2pt]{栈 $[7,3]$} (q1);
  \draw[credit] (L1) to[bend right=18] (q2);
  \draw[credit] (L2) to[bend right=18]
    node[lab,right=2pt]{栈 $[7,5]$} (q3);
  % 上行：request 按栈逐级聚合并累计贡献数
  \draw[up] (q1) to[bend right=18] (L1);
  \draw[up] (q2) to[bend right=18] (L1);
  \draw[up] (L1) to[bend right=18] node[lab,right=2pt]{$\times2$} (SP);
  \draw[up] (q3) to[bend right=18] (L2);
  \draw[up] (L2) to[bend right=18] node[lab,left=2pt]{$\times1$} (SP);
  \draw[up] (SP) to[bend right=18]
    node[lab,right=2pt]{$\times3$，完成样本} (R);
  % 图例
  \draw[credit] (-3.30,-0.55) -- (-2.85,-0.55);
  \node[lab,anchor=west] at (-2.80,-0.55) {credit 下行：分配 slot、压位置栈};
  \draw[up] ( 0.55,-0.55) -- ( 1.00,-0.55);
  \node[lab,anchor=west] at ( 1.05,-0.55) {request 上行：按栈聚合};
\end{tikzpicture}
\par\smallskip{\small (a) 一份 normal credit 的组播闭环}
\end{minipage}\hfill
\begin{minipage}[b]{0.40\linewidth}
\centering
\begin{tikzpicture}[x=1cm,y=1cm,
  box/.style={draw,rounded corners=1.5pt,minimum height=0.5cm,
    inner sep=3pt,font=\scriptsize},
  ep/.style={box,fill=ArborBlue!12},
  sw/.style={box,fill=ArborLightGray!60},
  t1/.style={ArborBlue,semithick},
  t2/.style={ArborOrange,semithick,dashed},
  lab/.style={font=\tiny,inner sep=1pt}]
  \node[sw] (SA) at (-0.9,2.05) {$S_A$};
  \node[sw] (SB) at ( 0.9,2.05) {$S_B$};
  \node[ep] (Rb) at (-2.4,0.40) {R};
  \node[ep] (b1) at (-0.8,0.40) {$q_1$};
  \node[ep] (b2) at ( 0.8,0.40) {$q_2$};
  \node[ep] (b3) at ( 2.4,0.40) {$q_3$};
  \foreach \h in {Rb,b1,b2,b3}{
    \draw[t2] (SB) -- (\h);
  }
  \foreach \h in {Rb,b1,b2,b3}{
    \draw[t1] (SA) -- (\h);
  }
  \node[lab,anchor=east,ArborBlue]   at (-1.38,2.42) {树 $s_1$};
  \node[lab,anchor=west,ArborOrange] at ( 1.38,2.42) {树 $s_2$};
  \node[lab,align=center] at (0,3.00)
    {offset：$k\to s_1$，\ $k{+}1\to s_2$，\ $k{+}2\to s_1$，$\cdots$};
  \node[lab,align=center] at (0,-0.30)
    {每棵树独立维护 $r_s$/$W_s$ gate；gate 关闭的树不取得新 offset};
\end{tikzpicture}
\par\smallskip{\small (b) 一个 channel 的多棵 subchannel 树}
\end{minipage}

\caption{多播-聚合流与组播 credit。(a) responder 为 packet offset $k$
在树 $s_1$ 上组播一份 credit；credit 副本逐级推进 allocator、取得 slot
并压入位置栈，requester 凭 credit 发送 request，上行按栈逐级聚合并累计
贡献数，整组到齐（$\times3$）后 responder 获得闭环完成样本。(b) 一个
channel 预装多棵 subchannel 树，responder 按各树的 gate 分配后续 offset；
同一 offset 的全部输入使用同一棵树。}
\label{fig:design-overview}
\end{figure*}

\sysname 把 credit 附在 response 上组播下发：带 credit 的 response 授权
一个 offset 和 subchannel，下行副本沿各自分支取得位置栈。requester 只有
消费该 credit 后才发送对应 request，并带回本分支的位置与逐级 fan-in。
上行聚合结果返回 responder 后，response 可以同时交付结果并授权新的
offset。注册闭环负责 message readiness，数据闭环负责 packet 许可与完成
（\cref{fig:design-overview}(a)）。

\parab{注册闭环.}
responder 只有在全组安装一条消息的本地状态后，才能为该消息发放 normal
credit。requester 提交消息后，在各 subchannel 上周期发送轻量
\texttt{REGISTER}；responder 为每个 message/subchannel 收集 requester
readiness，并在本地消息也已提交后组播 \texttt{REGISTER\_ACK}。周期发送
容忍控制包丢失，确认后 requester 停止对应信标。注册只确认 readiness，
不授予 request 发送。\sysname 只让正常路径的包经过加速器，其余包由网络
直接转发；拥塞控制也不采集注册阶段的样本。后续消息可以在当前消息传输期间
异步注册，因此注册闭环不要求 channel 在消息边界停顿。

\parab{Credit 驱动的数据闭环.}
responder 按每个 subchannel 的速率门和窗口门发放 credit。Credit 通过
response packet 组播给 requester group，使整组获得同一 offset 和
subchannel 的后续许可；response 副本沿各自分支携带聚合器位置，同一聚合
子树内的对应位置一致。readiness 完成后的数据 request 必须凭 credit 发送；
同一个 response 可以返回一个 packet 的数据或完成信息，并授权另一个
packet 的 request。负载均衡、拥塞控制与聚合器复用都由该 response/credit
闭环驱动。

\section{Channel 与 Subchannel}
\label{ssec:channel-subchannel}

channel 是多播-聚合流的协议表示：一个 channel 固定 responder 和
requester group，并依次承载多个 message。subchannel 是该 channel 在具体
拓扑上的一棵聚合树；每个 channel 独立定义自己的 subchannel 集合，不同
集合可以共享物理链路（\cref{fig:design-overview}(b)）。不同 packet
offset 可以选择不同 subchannel；同一 offset 的全部输入必须使用 credit
指定的同一棵树。

\Cref{tab:collective-mapping} 通过 channel 集合和用户 payload 方向表达
六种原语。``全部 rank''表示每个 rank 对应的 channel 都承载一个数据切片；
``root''只使用 root 对应的 channel。表中的有/无仅描述应用数据方向：
\texttt{Reduce} 和 \texttt{ReduceScatter} 的 response 仍携带 completion
与 credit，\texttt{Barrier} 的两个方向仍使用控制 packet。所有原语复用
相同的 credit、选树、聚合和恢复规则。

\begin{table}[t]
\centering
\caption{\sysname 在同一传输层数据通路上表达六种集合通信原语，差异仅在于
channel 集合的选择与 request/response 中用户 payload 方向的配置。}
\label{tab:collective-mapping}
\begin{tabular}{lccc}
\toprule
原语 & Channel & Request payload & Response payload \\
\midrule
AllReduce & 全部 rank & 有 & 有 \\
Reduce & root & 有 & 无 \\
Broadcast & root & 无 & 有 \\
AllGather & 全部 rank & 无 & 有 \\
ReduceScatter & 全部 rank & 有 & 无 \\
Barrier & 全部 rank & 无 & 无 \\
\bottomrule
\end{tabular}
\end{table}

\parab{Credit Context.}
在一个 channel 中，一份 credit 以 \texttt{(channel\_id, packet\_offset,
subchannel\_id)} 定义一个 \emph{Credit Context}。\texttt{packet\_offset}
标识本轮被授权的 packet，\texttt{subchannel\_id} 指定对应的聚合树。消费
同一份 credit 的 requester 继承同一 Credit Context；responder 按
\texttt{(channel\_id, packet\_offset)} 提交完成结果。

\section{聚合器位置反馈}
\label{ssec:agg-feedback}

聚合树上的每一级加速器维护一个 allocator 和一个共享 slot 池；同一加速器
的所有端口共用该池。Normal credit 沿下行路径经过各级加速器，每一级推进
allocator，将取得的 slot 的 owner 设为该 credit 的 Credit Context，并把
slot 编号压入位置栈。requester 因而只会收到已经取得沿途全部聚合位置的
normal credit。Repair trigger 不分配 slot。

\parab{Request 按位置栈逐级聚合.}
requester 把 credit 返回的位置栈和逐级 fan-in 写入 request。若本级
fan-in 为 1，加速器弹栈直转；否则按栈顶定位 slot，并验证 request 的
Credit Context 与 slot owner。首个需要合并的输入绑定 fan-in、算子和
payload shape；数据面同时检查输入是否属于当前支持的 data-type profile。
后续输入通过相同参数检查后累加贡献数和 payload。单个 requester 输入贡献
1，下层聚合结果贡献其携带的计数。贡献数达到本级 fan-in 后，加速器发出
aggregated packet、弹出本级位置，并由上一级继续处理。\texttt{Broadcast}
和 \texttt{AllGather} 使用同一计数路径合并空 request。

所有聚合参数均随 request 携带，加速器无需预装通信组、rank 或源地址表。
requester 对同一 Credit Context 至多发送一份 normal request；repair
request 直接旁路。因此，slot 保存 owner、聚合元数据、贡献数和累加结果，
但不保存 requester 身份或 per-requester bitmap。成员确认和 repair 去重由
responder 完成。

\parab{Owner 失配终止正常聚合.}
owner 失配表示 credit 许诺的 slot 已被新的 Credit Context 覆盖。加速器
丢弃 payload，并把 packet 改写为 header-only 的 \texttt{AGG\_MISS}；该
通知不再进入后续各级的正常聚合。responder 收到通知后将对应 credit 原子地
切换到 repair mode，并组播 repair trigger（\cref{ssec:reliability}）。
通知丢失时，原 credit 的超时器触发相同恢复过程。该规则同样处理单个输入和
下层 aggregated packet，不会把部分结果带入后续聚合级。

\parab{覆盖复用无需端到端释放.}
循环 allocator 通过覆盖复用 slot：normal credit 取得 slot 后，加速器
原子地安装新的 Credit Context owner 并重置逻辑聚合状态。旧 owner 的迟到
request 因 owner mismatch 产生 \texttt{AGG\_MISS}，responder 随后切入
repair。该规则把 slot 生命周期限制在本级：安全时刻由本级结果读出或
responder 的 repair cutover 确定。allocator 不等待端点或控制面发出显式
释放消息，也不维护分布式 aggregation-slot free list。加速器仅在把本级
结果复制到独立的完成缓冲区期间短暂阻止覆盖对应 slot；由此，slot 回收从
端到端协议状态变为加速器内部的局部覆盖操作。每一级独立维护循环指针和
owner；插入新的聚合层级只扩展 credit 的位置栈，不增加端点对该级 slot 的
分配或释放状态。\Cref{sssec:slot-bound} 给出不提前覆盖正常 Context 的
循环深度；深度不足时，被提前覆盖的 Context 转入 repair，旧部分结果仍
不能进入新 Context。

\section{Credit Context 与多路径选择}
\label{ssec:response-credit}
\label{ssec:packet-binding}

\parab{Credit Context 唯一性.}
多个 subchannel 的 response/credit 并行返回时，不同 subchannel 的 credit
在不同 requester 上可能以不同相对顺序到达。若 requester 按本地``下一个
未发送 packet''消费每个到达的 credit，同一份组播 credit 就可能在不同
requester 上绑定到不同 packet offset，使同一轮 fan-in 混入不同 packet 的
输入。\sysname 因此要求一份逻辑 credit 在所有 requester 上导出同一个
Credit Context，把聚合上下文一致性（\cref{ssec:group-consistency}）落实
为可检查的协议不变量。设组播 credit $C$ 逻辑上对应的 \texttt{channel\_id}、
\texttt{credit\_offset} 和 \texttt{subchannel\_id} 分别为 $c(C)$、
$k_c(C)$ 和 $s(C)$；requester $q \in Q$ 用 $C$ 实际发送的 request 继承的
三元组为 $(c_q(C), \mathrm{offset}_q(C), s_q(C))$。本文称如下条件为
\emph{Credit Context 唯一性}：
\begin{equation}
    \forall q \in Q,\quad
    (c_q(C), \mathrm{offset}_q(C), s_q(C))
    =
    (c(C), k_c(C), s(C)).
    \label{eq:credit-context-uniqueness}
\end{equation}

Credit Context 唯一性是组播 credit 的协议不变量。credit 确定 channel、
packet offset 和 subchannel；requester 据此选择 packet 和聚合树，并带回
沿途取得的位置栈（\cref{ssec:agg-feedback}）。未收到该 credit 的
requester 不会为对应 offset 发送 request。缺失输入和 response 由
\cref{ssec:reliability} 的恢复机制处理。

\parab{Payload/Credit offset 绑定.}
一个 response packet 可以返回 packet $p$ 的 payload 或完成信息，并用
credit 授权 packet $k$ 的 request。协议分别使用 \texttt{payload\_offset}
和 \texttt{credit\_offset} 绑定这两项语义；二者可以相同，也可以不同。
\texttt{Broadcast} 和 \texttt{AllGather} 主要使用前者，\texttt{Reduce}
和 \texttt{ReduceScatter} 主要使用后者，\texttt{AllReduce} 同时使用二者。
Credit Context 唯一性约束 credit 授权的 offset，不受 response payload
顺序影响。

\parab{反馈驱动的逐-offset 选树.}
responder 以 packet offset 为单位在多个预装 subchannel 间发放 credit。
每个 subchannel 都有速率门 $r_s$ 和窗口门 $W_s$；只有发放间隔到期且
未完成 credit 数小于 $W_s$ 时，该 subchannel 才能接收新的 packet
offset。responder 在当前允许发放的 subchannel 间分配后续 offset，拥塞或
恢复增多的 subchannel 因速率门收紧或窗口耗尽而获得更少 offset。已经发出
的 Credit Context 始终绑定原 subchannel，不在途换树。

\section{端到端拥塞控制}
\label{ssec:cc}

\parab{端点通过速率门和窗口门控制注入.}
\sysname 不依赖运行时集中式控制器：普通交换机的转发与组播表项在建组前
预装，responder 在本地更新每个 subchannel 的 $r_s$ 和 $W_s$，聚合加速器
不参与拥塞控制决策。requester 只凭 credit 发送，因此 responder 可以用
两道门控制每个 subchannel：速率门要求相邻 credit 的发放间隔不小于
$P/r_s$，窗口门要求未完成 credit 数小于 $W_s$，其中 $P$ 是一份 credit
授权的 packet 字节数，$r_s$ 和下式的 $x_s$ 均以 byte/s 计。令 $T_s$ 为该
subchannel 的稳态闭环完成时延，则 requester 分支上的稳态注入速率满足
\begin{equation}
    x_s \le \min\left\{r_s,\frac{W_sP}{T_s}\right\}.
    \label{eq:cc-two-gates}
\end{equation}
Rate-based、window-based 和 hybrid 控制律分别更新 $r_s$、$W_s$ 或二者。
小数窗口用 $W_s=1$ 加 pacing 表示。显式标记和基于完成时延的算法共用该
credit 发放接口，具体映射见 \cref{app:rate-modules}。

\parab{闭环完成时延覆盖当前最慢分支.}
responder 记录每份 credit 的发出时刻；当返回结果首次证明全组输入已经完成
时，两者之差给出其\emph{闭环完成时延}。该时间覆盖 credit 组播下行、
requester 消费、request 上行和 fan-in 聚合。令 $\mathcal{B}_s$ 表示
subchannel $s$ 上从 responder 经 requester 再返回的分支集合，$d_b$ 表示
分支 $b$ 对本轮完成的时延，则
\begin{equation}
    T_s=\max_{b\in\mathcal{B}_s} d_b.
    \label{eq:completion-delay}
\end{equation}
requester 发送路径停顿和网络排队都可以进入 $d_b$；消息提交偏斜由此前的
注册闭环吸收，不进入 normal credit 样本。只有决定最大值的分支直接改变
本轮完成时间。一份 credit、一次整组完成和一个 $T_s$ 使整组 requester
共享同一速率样本，无需沿途设备维护 per-flow 速率状态。

\parab{最大完成时延过滤非关键分支排队.}
考虑共享链路 $E$ 上增加的闭环排队时间 $q_E$。对一条经过 $E$ 的 channel
$i$，令 $S_i$ 为经过 $E$ 的分支中最大的无队列完成时延，$M_i$ 为其余分支
中的最大值；若所有分支都经过 $E$，定义 $M_i=-\infty$。于是
\begin{equation}
    T_i(q_E)=\max\{M_i,S_i+q_E\},\qquad
    \Delta_i=\max\{0,M_i-S_i\}.
    \label{eq:branch-slack}
\end{equation}
当 $q_E<\Delta_i$ 时，链路 $E$ 尚未决定该 channel 的完成时间，基于完成
时延的模块不会因这部分排队降低整组速率。此时降速不能直接缩短该 channel
的 fan-in 完成时间，反而会减少真正关键分支获得的流量。$q_E$ 达到
$\Delta_i$ 后，经过 $E$ 的分支成为当前最慢分支，后续排队完整进入 $T_i$。
因此，完成时延只让影响 collective 进度的分支直接参与速率调节。

\parab{共享链路上的公平性取决于分支余量.}
多个 channel 共享 $E$ 时，其 $\Delta_i$ 可以不同。若 $E$ 已是 channel
$j$ 的关键分支，却仍落在 channel $i$ 的非关键分支上，$j$ 会先观察到排队
并降速，$i$ 在 $q_E<\Delta_i$ 期间保持原速率。当 $i$ 已受其它分支限制且
不超过 $E$ 的剩余容量时，该行为避免无效退让；否则，非对称路径、
responder 放置和共享链路竞争会产生放置相关的公平性偏置。若所有竞争
channel 都以 $E$ 为关键分支，行为退化为所选控制律在共享瓶颈上的行为；
ECN 则可独立暴露任意分支上的队列。\Cref{ssec:eval-cc} 量化该吞吐与
公平性权衡。

\parab{参考状态与样本筛选.}
responder 在发放 credit 时保存所选控制律所需的参考状态，使返回样本按
发出时状态解释，避免用已经更新过的当前状态重复响应同一次拥塞。一份
normal credit 首次完成时产生一个样本；\texttt{AGG\_MISS}、repair、重放和
已经提交后的迟到输入不产生新样本，修复流量由 \cref{ssec:reliability} 的
独立预算限速。显式标记、rate-based、window-based 与 hybrid 模块所需信号
及其在该接口上的对应关系见\cref{app:rate-modules}。

\parab{显式 ECN 只需传播整组标记.}
基于 completion delay 的端侧算法只使用 responder 本地状态，不要求数据面
增加拥塞控制逻辑。可选 ECN 算法让 requester 把下行 credit 上的 CE 回显到
下一份 normal request，并在逐级聚合时对各输入的 CE 取 OR。responder 因而
获得覆盖上下行路径和全组分支的单个标记。该传播不引入 per-flow 拥塞状态，
控制律仍只在 responder 执行（\cref{app:rate-modules}）。

\section{可靠性与重传}
\label{ssec:reliability}

\sysname 将正常路径的纯计数与恢复路径的身份去重分开。requester 在正常
路径保证单注入，加速器按贡献数完成聚合；异常发生后，responder 将整个
Credit Context 切换到 repair，并用 requester 位图完成去重和回退聚合。

\parab{Normal request 单注入.}
requester 对每个 Credit Context 至多发出一份 normal request。同一份
credit 在 requester 本地只消费一次；之后对该 packet offset 的重发只由
repair trigger 触发，并携带 repair 标志。repair request 旁路 slot。
Normal request 在上行路径不被复制，因此同一 Credit Context 可以按贡献数
聚合；repair 去重和幂等提交由 responder 位图处理。

\parab{Responder 单点提交与重传.}
responder 为每份 normal credit 保留未完成记录，并采用与 ATP 类似的单次
提交语义~\cite{atp}。正常聚合结果只有在贡献数等于 requester group 规模时
才能提交；提交后缓存的 response 不再被迟到输入改写。responder 是唯一的
超时检测与恢复发起者，避免多个 requester 对同一份组播 credit 独立产生
NACK。超时或 \texttt{AGG\_MISS} 将未完成 credit 原子地切换到 repair
mode；此后迟到的正常聚合结果不能提交，该 Credit Context 只由 repair 路径
完成。

\parab{组播 repair trigger.}
定时器超时或收到交换机的失配通知（\cref{ssec:agg-feedback}）后，
responder 组播一份 repair trigger，携带 \texttt{(channel\_id,
packet\_offset, subchannel\_id)}。该控制消息标识恢复范围，requester 据此
沿原 subchannel 发送旁路正常聚合的 repair request。repair trigger 必须
组播：responder 只能观察最终到达的 aggregated packet、\texttt{AGG\_MISS}
和 repair request，看不到交换机 slot 中已经累加的部分状态，也无法可靠
判断哪一个 requester 的输入缺失。Repair request 不再进入原 slot，而是
直达 responder。因此恢复边界是整个 requester group：所有 requester 重新
提交，responder 再用位图收齐并去重。Repair 流量使用独立保守预算，且不
产生正常完成样本（\cref{ssec:cc}）。

\parab{Response 重放不创建第二份 normal credit.}
Response 重放不会使同一 requester 为一个 Credit Context 发送第二份
normal request：未完成 credit 以 repair trigger 重发；缓存的尾部
response 保留 normal credit 时，requester 按本地消费记录去重。已经提交的
结果直接从缓存 response 重放，不重新执行聚合。重试耗尽后，该 collective
失败，且不再复用其未确认的 message identifier。

\parab{END 确认消息完成并回收状态.}
消息末尾若干 packet 的 response 之后不再有新的 credit，其交付无法由后续
credit 的消费隐式确认。全部 credit 完成后，responder 组播 \texttt{END}；
requester 在本地消息完成后回复 \texttt{END\_ACK}。确认缺失时，responder
重放尚未被后续 credit 确认的尾部 response，并重新发送 \texttt{END}。握手
完成前，responder 保留相应 response 和消息状态；握手只延迟消息标识符
复用，不阻止已经注册的后续消息使用新的 packet offset。

\section{共享循环 allocator 的容量}
\label{ssec:bounded-state}
\label{sssec:slot-bound}

每台交换机的全部端口和 job 共用一个带 owner 的 slot 池，每个 normal
Credit Context 都推进同一个循环 allocator。本节要回答的问题可以说得很
直白：池子要配多少个 slot，才能保证 allocator 转一圈回到某个地址时，
上面的旧聚合一定已经结束？答案分两步：先给出一个精确的充要条件，再把
它换算成只依赖端口线速的部署上界。

\parab{循环复用的基本结构.}
设池中有 $A_e$ 个 slot，第 $n$ 次 allocation 使用地址 $n\bmod A_e$。
于是同一个地址恰好每隔 $A_e$ 次 allocation 被复用一次：某个 Context
在第 $n$ 次 allocation 拿到 slot 后，下一个使用同址的正是第 $n+A_e$
次 allocation。``复用是否安全''因此只取决于一个数字——从拿到 slot 到
它可以被安全覆盖，allocator 中间又推进了多少次。只要这个数字不超过
$A_e$，轮回到同址时旧 Context 必然已经结束。下面把``可以被安全覆盖''
和``推进了多少次''说精确。

\parab{复用保护区间.}
覆盖一个 slot 唯一的危险，是旧 Context 之后仍可能在本级提交 normal
结果。两类事件都会永久排除这种可能：本级聚合结果已经读出，或
responder 已完成 repair cutover——此后迟到的 normal 结果不再能提交
（\cref{ssec:reliability}）。对经过交换机 $e$ 的 normal Credit
Context $x$，记 $t_a(e,x)$ 为绑定 slot 的时刻，$t_f(e,x)$ 为上述两类
事件中先发生的时刻；$t_f$ 是 $x$ 的\emph{可安全覆盖时刻}，而非协议中
的显式释放事件。半开区间 $[t_a(e,x),t_f(e,x))$ 称为 $x$ 的\emph{复用
保护区间}，取 $R_e>0$ 为其长度的上界。

\parab{精确最小深度.}
令 $N_e(I)$ 表示事件区间 $I$ 内的 normal advance 数，对保护区间从
$x$ 自身的 allocation（含）计到 $t_f$（不含）。所需深度由最坏的
Context 决定：
\begin{equation}
 D_e=\sup_x N_e([t_a(e,x),t_f(e,x))).
 \label{eq:allocator-exact-depth}
\end{equation}
用一句话读这个定义：$D_e$ 是``任何一个 Context 在自己还不能被覆盖的
期间，整台交换机最多又发生了多少次分配''。计数覆盖全部 normal
advance——fan-in 为 1、弹栈直转因而自身可立即覆盖的 Context 也占一次
分配，同样推进循环指针。$D_e$ 只关于全部 job 合并后的同一条分配序列，
池容量因此不需要按 job 或 rank 数再乘静态因子。同一时刻的 $t_f$ 与 allocation 按先释放、后
分配排序；完整记账约定与实现无法保证该顺序时的处理见
\cref{app:allocator-port-bound}。

\begin{theorem}
\label{thm:allocator-exact}
第 $n$ 次 allocation 使用地址 $n\bmod A_e$ 时，normal path 不发生提前
覆盖当且仅当
\begin{equation}
 A_e\geq D_e .
 \label{eq:allocator-exact-condition}
\end{equation}
\end{theorem}

\noindent\emph{证明.}
充分性：任取 Context $x$，设它由第 $n$ 次 allocation 绑定，则复用
同址的是第 $n+A_e$ 次。由 $A_e\geq D_e$，$x$ 的保护区间内至多有
$A_e$ 次 advance（含 $x$ 自身），即区间内的 allocation 至多编号到
$n+A_e-1$；第 $n+A_e$ 次因此不早于 $t_f$——覆盖发生时 $x$ 已可安全
覆盖。必要性：若 $A_e<D_e$，按 \cref{eq:allocator-exact-depth} 存在
Context $x$，其保护区间内的 advance 数超过 $A_e$，于是第 $n+A_e$ 次
allocation 落在 $t_f$ 之前——同一地址在 $x$ 可安全覆盖前就被覆盖。
\Cref{fig:allocator-timeline} 画出这两种情形。\hfill$\square$

\begin{figure}[t]
\centering
\begin{tikzpicture}[font=\scriptsize]
\draw[->,black!70] (-0.2,0) -- (8.6,0);
\fill[ArborBlue,fill opacity=0.12] (0.6,0) rectangle (5.4,0.72);
\draw[ArborBlue,thick] (0.6,-0.14) -- (0.6,0.72);
\draw[ArborBlue,thick,dashed] (5.4,-0.14) -- (5.4,0.72);
\node[ArborBlue!70!black,anchor=south] at (2.8,0.76)
  {复用保护区间 $[t_a,t_f)$，长度 $\leq R_e$};
\node[anchor=north,align=center] at (0.8,-0.22)
  {$t_a$：第 $n$ 次 allocation\\绑定地址 $n\bmod A_e$};
\node[anchor=north,align=center] at (5.6,-0.22) {$t_f$：可安全覆盖};
\foreach \x in {1.3,2.0,2.7,3.4,4.1} \draw[black!55] (\x,-0.09) -- (\x,0.09);
\draw[black!55] (7.1,-0.09) -- (7.1,0.09);
\node[black!55,anchor=south] at (2.35,0.10) {后续 normal advance};
\draw[ArborOrange,thick,->] (4.8,1.46) -- (4.8,0.12);
\node[ArborOrange!85!black,anchor=south,align=center] at (4.3,1.48)
  {第 $n{+}A_e$ 次（$A_e<D_e$）\\提前覆盖};
\draw[ArborGreen!80!black,thick,->] (7.1,1.46) -- (7.1,0.12);
\node[ArborGreen!55!black,anchor=south,align=center] at (7.3,1.48)
  {第 $n{+}A_e$ 次（$A_e\geq D_e$）\\安全复用};
\end{tikzpicture}
\caption{循环复用的时间线。地址 $n\bmod A_e$ 在第 $n{+}A_e$ 次
allocation 时被再次使用：保护区间内的 advance 数不超过 $A_e$（即
$A_e\geq D_e$）时，同址复用必然落在可安全覆盖时刻 $t_f$ 之后；否则第
$n{+}A_e$ 次分配落入区间内，构成提前覆盖。}
\label{fig:allocator-timeline}
\end{figure}

容量不足不改变结果正确性。新 allocation 原子地安装新的 owner 并重置
逻辑聚合状态；覆盖后到达的旧 request 无法通过 owner 检查，并被改写为
\texttt{AGG\_MISS}。若不再有旧 request 到达，responder timeout 仍会
切入 repair。Normal request 单注入又阻止旧 Context 开始第二轮聚合。
欠配因而增加 repair，但不会把旧部分结果提交给新的 Context。

\parab{端口速率把 $D_e$ 换算为部署深度.}
$D_e$ 是协议量，部署者需要能直接从链路参数算出的上界。物理直觉是：
每一次 advance 最终都对应一份至少 $P$ bytes 的 payload packet 经过
某个端口，而端口串行化限制了任何时间窗内能通过的这类包数——一个
100-Gbps 端口发送一个 8-KiB 包约需 0.66~\textmu s，长度 $R_e$ 的窗口
内它至多承载约 $R_e/(0.66~\mu\mathrm{s})$ 次由它见证的推进。形式化地，考虑
每个 advance 都对应一个正常完成的 fixed-$P$ data context；若每个
context 都能与 allocation 后 $R_e$ 内一份至少携带 $P$ bytes 的
payload packet 一一对应，令 $\mathcal P_e$ 为这些 packet 实际经过的
有向端口集合，$C_p$ 为端口 $p$ 的线速，则
\begin{equation}
 D_e\leq
 \sum_{p\in\mathcal P_e}
 \left\lceil\frac{2C_pR_e}{8P}\right\rceil .
 \label{eq:allocator-directed-port}
\end{equation}
若推进 allocator 的 packet 本身就携带至少 $P$ bytes，配对区间从
$2R_e$ 收紧为 $R_e$：
\begin{equation}
 D_e\leq
 \sum_{p\in\mathcal P_e}
 \left\lceil\frac{C_pR_e}{8P}\right\rceil .
 \label{eq:allocator-payload-carrier-bound}
\end{equation}
两条界均按实际推进同一个共享 allocator 的端口求和。payload 配对与
两倍区间的证明，以及面向提供 advance 到达包络保证的实现、不依赖
payload 配对的另一条充分条件，见 \cref{app:allocator-port-bound}。

在 payload 一一配对前提仍成立时，若部署通过 admission、shaping 和
normal-to-repair deadline 保证保护时间不超过 $R_{\max}$，就可以在
上述容量式中直接使用 $R_{\max}$。沿用上面的直觉：100-Gbps 端口、
$P=8192$ bytes、$R_{\max}=100$~\textmu s 时，
\cref{eq:allocator-payload-carrier-bound} 给出 153 个 slot，即约
1.20~MiB payload memory，另加 metadata。部署使用契约保证的
$R_{\max}$，而非运行 trace 中观察到的最大值。

由循环复用的基本结构，任意连续 $A_e$ 次 allocation 恰好访问每个地址
一次；$A_e\geq D_e$ 时，全部已配置 slot 都是无冲突容量，不受哈希冲突
或 per-job 分区碎片影响。\Cref{ssec:eval-state} 在共用端口和端口分散
两种放置中核对实测 $D_e$ 与 \cref{eq:allocator-payload-carrier-bound}。
